\documentclass{beamer}
\usetheme{Madrid}
\usecolortheme{default}
\usepackage[utf8]{inputenc}
\usepackage{amsmath}
\usepackage{amsfonts}
\usepackage{amssymb}
\usepackage{tikz}

\title{CRISC Exam - Hard Questions with Explanations}
\subtitle{From   CRISC GPTs}
\author{Compiled from ISACA CRISC GPT APIs}
\date{\today}

\begin{document}
	
	\begin{frame}
		\titlepage
	\end{frame}
	
	\begin{frame}{Outline}
		\tableofcontents
	\end{frame}
	
	% ================ SECTION 1: DOMAIN 1 - IT RISK IDENTIFICATION ================
	
	\section{Domain 1: IT Risk Identification}
	
	\begin{frame}{Q1: Resilient Business Processes}
		\begin{block}{Question}
			Which business requirement BEST relates to the need for resilient business and information systems processes?
		\end{block}
		\begin{enumerate}[(A)]
			\item Effectiveness
			\item Confidentiality
			\item Integrity
			\item Availability
		\end{enumerate}
	\end{frame}
	
	\begin{frame}{Q1: Explanation}
		\begin{block}{Correct Answer: D}
			Availability relates to information being available when required by the business process now and in the future. Resilience is the ability to provide and maintain an acceptable level of service during disasters or when facing operational challenges.
		\end{block}
		\begin{block}{Option Analysis}
			\begin{itemize}
				\item \textbf{A - Incorrect:} Effectiveness deals with information being relevant and delivered timely
				\item \textbf{B - Incorrect:} Confidentiality deals with protection from unauthorized disclosure
				\item \textbf{C - Incorrect:} Integrity relates to accuracy and completeness
				\item \textbf{D - Correct:} Availability is directly linked to resilience
			\end{itemize}
		\end{block}
	\end{frame}
	
	\begin{frame}{Q2: Value of Risk Register}
		\begin{block}{Question}
			Which statement BEST describes the value of a risk register?
		\end{block}
		\begin{enumerate}[(A)]
			\item It captures the risk inventory
			\item It drives the risk response plan
			\item It is a risk reporting tool
			\item It lists internal risk and external risk
		\end{enumerate}
	\end{frame}
	
	\begin{frame}{Q2: Explanation}
		\begin{block}{Correct Answer: B}
			Risk registers serve as the main reference for all risk-related information, supporting risk-related decisions such as risk response activities and their prioritization.
		\end{block}
		\begin{block}{Option Analysis}
			\begin{itemize}
				\item \textbf{A - Incorrect:} A risk register provides detailed information on each risk, not just inventory
				\item \textbf{B - Correct:} Drives risk response decisions
				\item \textbf{C - Incorrect:} Register data generate reports but is not itself a reporting tool
				\item \textbf{D - Incorrect:} Tracks all risk but this is not its primary value
			\end{itemize}
		\end{block}
	\end{frame}
	
	\begin{frame}{Q3: New Law Impact on HR System}
		\begin{block}{Question}
			Shortly after annual policy review, a new law may affect security requirements for HR system. Risk practitioner should:
		\end{block}
		\begin{enumerate}[(A)]
			\item Analyze how the law may affect the enterprise
			\item Ensure adjustments during next review cycle
			\item Initiate an ad hoc revision of corporate policy
			\item Notify system custodian to implement changes
		\end{enumerate}
	\end{frame}
	
	\begin{frame}{Q3: Explanation}
		\begin{block}{Correct Answer: A}
			Assessing how the law may affect the enterprise is the best course of action. The analysis must also determine whether existing controls already address the new requirements.
		\end{block}
		\begin{block}{Option Analysis}
			\begin{itemize}
				\item \textbf{A - Correct:} Analysis is the first step
				\item \textbf{B - Incorrect:} Significant changes should trigger ad hoc reassessment
				\item \textbf{C - Incorrect:} Policy should be developed systematically
				\item \textbf{D - Incorrect:} Changes need approval by process owner
			\end{itemize}
		\end{block}
	\end{frame}
	
	\begin{frame}{Q4: Weather Forecast System Priority}
		\begin{block}{Question}
			An information system that processes weather forecasts for public consumption is MOST likely to place its highest priority on:
		\end{block}
		\begin{enumerate}[(A)]
			\item Nonrepudiation
			\item Confidentiality
			\item Integrity
			\item Availability
		\end{enumerate}
	\end{frame}
	
	\begin{frame}{Q4: Explanation}
		\begin{block}{Correct Answer: C}
			A system that delivers weather forecasts is likely to place its highest priority on the integrity of the data. The risk practitioner should keep in mind that whether a forecast turns out to be accurate in its prediction is distinct from whether the data was accurately represented.
		\end{block}
		\begin{block}{Option Analysis}
			\begin{itemize}
				\item \textbf{A - Incorrect:} Nonrepudiation is unlikely to be important for weather forecasts
				\item \textbf{B - Incorrect:} Keeping data confidential would be at odds with public use
				\item \textbf{C - Correct:} Integrity of data is paramount
				\item \textbf{D - Incorrect:} Availability is important but integrity is higher priority
			\end{itemize}
		\end{block}
	\end{frame}
	
	\begin{frame}{Q5: Best Risk Management View}
		\begin{block}{Question}
			Which provides the BEST view of risk management?
		\end{block}
		\begin{enumerate}[(A)]
			\item An interdisciplinary team
			\item A third-party risk assessment
			\item IT department risk assessment
			\item Internal compliance department
		\end{enumerate}
	\end{frame}
	
	\begin{frame}{Q5: Explanation}
		\begin{block}{Correct Answer: A}
			Having an interdisciplinary team contribute to risk management ensures that all areas are adequately considered and included in the risk assessment processes to support an enterprise view of risk.
		\end{block}
		\begin{block}{Option Analysis}
			\begin{itemize}
				\item \textbf{A - Correct:} Interdisciplinary team provides enterprise-wide view
				\item \textbf{B - Incorrect:} Without internal knowledge, assessment adequacy is difficult
				\item \textbf{C - Incorrect:} IT department unlikely to reflect entire enterprise
				\item \textbf{D - Incorrect:} Compliance department focuses on regulatory implications
			\end{itemize}
		\end{block}
	\end{frame}
	
	\begin{frame}{Q6: IT Risk Awareness Program}
		\begin{block}{Question}
			Which is PRIMARY consideration when developing an IT risk awareness program?
		\end{block}
		\begin{enumerate}[(A)]
			\item Why technology risk is owned by IT
			\item How technology risk can impact each attendee's area of business
			\item How business process owners can transfer technology risk
			\item Why technology risk is more difficult to manage
		\end{enumerate}
	\end{frame}
	
	\begin{frame}{Q6: Explanation}
		\begin{block}{Correct Answer: B}
			Stakeholders must understand how the IT-related risk impacts the overall business.
		\end{block}
		\begin{block}{Option Analysis}
			\begin{itemize}
				\item \textbf{A - Incorrect:} IT does not own technology risk; business owns it
				\item \textbf{B - Correct:} Business impact is the primary consideration
				\item \textbf{C - Incorrect:} Transferring risk is part of risk response, not awareness
				\item \textbf{D - Incorrect:} Difficulty of management is secondary to business impact
			\end{itemize}
		\end{block}
	\end{frame}
	
	\begin{frame}{Q7: Risk Appetite Alignment}
		\begin{block}{Question}
			It is MOST important that risk appetite is aligned with business objectives to ensure that:
		\end{block}
		\begin{enumerate}[(A)]
			\item Resources are directed toward areas of low risk tolerance
			\item Major risk is identified and eliminated
			\item IT and business goals are aligned
			\item Risk strategy is adequately communicated
		\end{enumerate}
	\end{frame}
	
	\begin{frame}{Q7: Explanation}
		\begin{block}{Correct Answer: A}
			Risk appetite is the amount of risk that an enterprise is willing to take on in pursuit of value. Aligning it with business objectives allows an enterprise to evaluate and deploy valuable resources toward those objectives where the risk tolerance (for loss) is low.
		\end{block}
		\begin{block}{Option Analysis}
			\begin{itemize}
				\item \textbf{A - Correct:} Resources directed to low tolerance areas
				\item \textbf{B - Incorrect:} Risk elimination is rarely cost-effective
				\item \textbf{C - Incorrect:} Alignment affects more than IT and business
				\item \textbf{D - Incorrect:} Communication does not depend on alignment
			\end{itemize}
		\end{block}
	\end{frame}
	
	\begin{frame}{Q8: Weak Passwords and Unprotected Transmission}
		\begin{block}{Question}
			Weak passwords and transmission over unprotected communication lines are examples of:
		\end{block}
		\begin{enumerate}[(A)]
			\item Vulnerabilities
			\item Threats
			\item Probabilities
			\item Impacts
		\end{enumerate}
	\end{frame}
	
	\begin{frame}{Q8: Explanation}
		\begin{block}{Correct Answer: A}
			Vulnerabilities represent characteristics of information resources that may be exploited by a threat.
		\end{block}
		\begin{block}{Option Analysis}
			\begin{itemize}
				\item \textbf{A - Correct:} Weak passwords and unprotected transmission are vulnerabilities
				\item \textbf{B - Incorrect:} Threats are circumstances with potential to cause harm
				\item \textbf{C - Incorrect:} Probabilities represent likelihood of occurrence
				\item \textbf{D - Incorrect:} Impacts represent outcome of threat exploiting vulnerability
			\end{itemize}
		\end{block}
	\end{frame}
	
	% ================ SECTION 2: DOMAIN 2 - IT RISK ASSESSMENT ================
	
	\section{Domain 2: IT Risk Assessment}
	
	\begin{frame}{Q9: Quantitative Risk Analysis Drawback}
		\begin{block}{Question}
			MOST significant drawback of using quantitative risk analysis instead of qualitative risk analysis is:
		\end{block}
		\begin{enumerate}[(A)]
			\item Lower objectivity
			\item Greater reliance on expertise
			\item Less management buy-in
			\item Higher cost
		\end{enumerate}
	\end{frame}
	
	\begin{frame}{Q9: Explanation}
		\begin{block}{Correct Answer: D}
			Quantitative risk analysis is generally more complex and, therefore, more costly than qualitative risk analysis.
		\end{block}
		\begin{block}{Option Analysis}
			\begin{itemize}
				\item \textbf{A - Incorrect:} Neither method is fully objective
				\item \textbf{B - Incorrect:} Both require personnel with good business understanding
				\item \textbf{C - Incorrect:} Quantitative generally has better buy-in
				\item \textbf{D - Correct:} Quantitative is more complex and costly
			\end{itemize}
		\end{block}
	\end{frame}
	
	\begin{frame}{Q10: Risk Scenario Analysis Purpose}
		\begin{block}{Question}
			Risk scenarios are analyzed to determine:
		\end{block}
		\begin{enumerate}[(A)]
			\item Strength of controls
			\item Likelihood and impact
			\item Current risk profile
			\item Scenario root cause
		\end{enumerate}
	\end{frame}
	
	\begin{frame}{Q10: Explanation}
		\begin{block}{Correct Answer: B}
			Risk scenarios are descriptions of events that can lead to a business impact and are evaluated to determine the likelihood and impact should the event occur.
		\end{block}
		\begin{block}{Option Analysis}
			\begin{itemize}
				\item \textbf{A - Incorrect:} Control strength is determined after controls are in place
				\item \textbf{B - Correct:} Likelihood and impact are the key assessments
				\item \textbf{C - Incorrect:} Risk profile is identification of current risk concerns
				\item \textbf{D - Incorrect:} Root cause is not part of risk scenario process
			\end{itemize}
		\end{block}
	\end{frame}
	
	\begin{frame}{Q11: Risk Ownership of Critical System}
		\begin{block}{Question}
			The risk to an information system that supports a critical business process is owned by:
		\end{block}
		\begin{enumerate}[(A)]
			\item IT director
			\item Senior management
			\item Risk management department
			\item System users
		\end{enumerate}
	\end{frame}
	
	\begin{frame}{Q11: Explanation}
		\begin{block}{Correct Answer: B}
			Senior management is responsible for the acceptance and mitigation of all risk.
		\end{block}
		\begin{block}{Option Analysis}
			\begin{itemize}
				\item \textbf{A - Incorrect:} IT director manages systems on behalf of business owners
				\item \textbf{B - Correct:} Senior management owns risk
				\item \textbf{C - Incorrect:} Risk management determines and reports on risk but does not own it
				\item \textbf{D - Incorrect:} System users use system properly but do not own risk
			\end{itemize}
		\end{block}
	\end{frame}
	
	\begin{frame}{Q12: Periodic Risk Assessment Reason}
		\begin{block}{Question}
			PRIMARY reason risk assessments should be repeated at regular intervals is:
		\end{block}
		\begin{enumerate}[(A)]
			\item Omissions in earlier assessments can be addressed
			\item Periodic assessments allow various methodologies
			\item Business threats are constantly changing
			\item They help raise risk awareness among staff
		\end{enumerate}
	\end{frame}
	
	\begin{frame}{Q12: Explanation}
		\begin{block}{Correct Answer: C}
			As business objectives and methods change, the nature and relevance of threats also change. This is the primary reason to conduct periodic risk assessments.
		\end{block}
		\begin{block}{Option Analysis}
			\begin{itemize}
				\item \textbf{A - Incorrect:} This is a benefit but not the primary reason
				\item \textbf{B - Incorrect:} Organizations strive to improve but this is not primary
				\item \textbf{C - Correct:} Threats constantly change
				\item \textbf{D - Incorrect:} Risk awareness is a minor benefit
			\end{itemize}
		\end{block}
	\end{frame}
	
	\begin{frame}{Q13: Measuring Risk Management Development}
		\begin{block}{Question}
			Which BEST assists in measuring existing level of development of risk management processes against desired state?
		\end{block}
		\begin{enumerate}[(A)]
			\item Capability maturity model (CMM)
			\item Risk management audit reports
			\item Balanced scorecard
			\item Enterprise security architecture
		\end{enumerate}
	\end{frame}
	
	\begin{frame}{Q13: Explanation}
		\begin{block}{Correct Answer: A}
			The capability maturity model (CMM) grades processes on a scale of 0 to 5, based on their maturity. It is commonly used by entities to measure their existing state and then to determine the desired one.
		\end{block}
		\begin{block}{Option Analysis}
			\begin{itemize}
				\item \textbf{A - Correct:} CMM measures current vs desired state
				\item \textbf{B - Incorrect:} Audit reports offer limited view
				\item \textbf{C - Incorrect:} Balanced scorecard measures implementation of strategy
				\item \textbf{D - Incorrect:} Enterprise security architecture explains business strategy view
			\end{itemize}
		\end{block}
	\end{frame}
	
	\begin{frame}{Q14: Identifying Control Deficiencies}
		\begin{block}{Question}
			Which BEST helps identify information systems control deficiencies?
		\end{block}
		\begin{enumerate}[(A)]
			\item Gap analysis
			\item Current IT risk profile
			\item IT controls framework
			\item Countermeasure analysis
		\end{enumerate}
	\end{frame}
	
	\begin{frame}{Q14: Explanation}
		\begin{block}{Correct Answer: A}
			Controls are deployed to achieve desired control objectives based on risk assessments and business requirements. The gap between desired control objectives and actual IS control design and operational effectiveness identifies control deficiencies.
		\end{block}
		\begin{block}{Option Analysis}
			\begin{itemize}
				\item \textbf{A - Correct:} Gap analysis identifies deficiencies
				\item \textbf{B - Incorrect:} Risk profile does not expose gaps
				\item \textbf{C - Incorrect:} Framework is generic without current state info
				\item \textbf{D - Incorrect:} Only identifies countermeasure deficiencies
			\end{itemize}
		\end{block}
	\end{frame}
	
	\begin{frame}{Q15: Statistical Methods Risk Analysis}
		\begin{block}{Question}
			Deriving likelihood and impact of risk scenarios through statistical methods is MOST LIKELY associated with which type of risk analysis?
		\end{block}
		\begin{enumerate}[(A)]
			\item Risk scenario
			\item Qualitative
			\item Quantitative
			\item Semiquantitative
		\end{enumerate}
	\end{frame}
	
	\begin{frame}{Q15: Explanation}
		\begin{block}{Correct Answer: C}
			The essence of quantitative risk assessment is to derive the likelihood and impact of risk scenarios based on statistical methods and data.
		\end{block}
		\begin{block}{Option Analysis}
			\begin{itemize}
				\item \textbf{A - Incorrect:} Scenario analysis includes several methods
				\item \textbf{B - Incorrect:} Qualitative uses experiential measures
				\item \textbf{C - Correct:} Quantitative uses statistical methods
				\item \textbf{D - Incorrect:} Semiquantitative assigns values to qualitative ratings
			\end{itemize}
		\end{block}
	\end{frame}
	
	\begin{frame}{Q16: Reviewing Risk Analysis Results}
		\begin{block}{Question}
			Which review is BEST suited for review of IT risk analysis results before sending to management?
		\end{block}
		\begin{enumerate}[(A)]
			\item Internal audit review
			\item Peer review
			\item Compliance review
			\item Risk policy review
		\end{enumerate}
	\end{frame}
	
	\begin{frame}{Q16: Explanation}
		\begin{block}{Correct Answer: B}
			It is effective, efficient and good practice to perform a peer review of IT risk analysis results before sending them to management.
		\end{block}
		\begin{block}{Option Analysis}
			\begin{itemize}
				\item \textbf{A - Incorrect:} Audit is independent assurance
				\item \textbf{B - Correct:} Peer review is good practice
				\item \textbf{C - Incorrect:} Compliance reviews measure conformance
				\item \textbf{D - Incorrect:} Policy review may change content but isn't a review method
			\end{itemize}
		\end{block}
	\end{frame}
	
	% ================ SECTION 3: DOMAIN 3 - RISK RESPONSE ================
	
	\section{Domain 3: Risk Response and Mitigation}
	
	\begin{frame}{Q17: Risk Cannot Be Sufficiently Mitigated}
		\begin{block}{Question}
			When risk cannot be sufficiently mitigated through manual or automatic controls, which option BEST protects from potential financial impact?
		\end{block}
		\begin{enumerate}[(A)]
			\item Insuring against the risk
			\item Updating IT risk register
			\item Improving staff training
			\item Outsourcing related process
		\end{enumerate}
	\end{frame}
	
	\begin{frame}{Q17: Explanation}
		\begin{block}{Correct Answer: A}
			An insurance policy can compensate the enterprise monetarily for the impact of the risk by transferring the risk to the insurance company.
		\end{block}
		\begin{block}{Option Analysis}
			\begin{itemize}
				\item \textbf{A - Correct:} Insurance transfers financial risk
				\item \textbf{B - Incorrect:} Updating register doesn't change risk
				\item \textbf{C - Incorrect:} Staff training may reduce but not completely mitigate
				\item \textbf{D - Incorrect:} Outsourcing does not necessarily remove risk
			\end{itemize}
		\end{block}
	\end{frame}
	
	\begin{frame}{Q18: Risk Mitigation Must Reduce}
		\begin{block}{Question}
			To be effective, risk mitigation MUST reduce the:
		\end{block}
		\begin{enumerate}[(A)]
			\item Residual risk
			\item Inherent risk
			\item Frequency of a threat
			\item Impact of a threat
		\end{enumerate}
	\end{frame}
	
	\begin{frame}{Q18: Explanation}
		\begin{block}{Correct Answer: A}
			The objective of risk reduction is to reduce the residual risk to levels below the enterprise's risk tolerance level.
		\end{block}
		\begin{block}{Option Analysis}
			\begin{itemize}
				\item \textbf{A - Correct:} Residual risk must be reduced to acceptable level
				\item \textbf{B - Incorrect:} Reduction of inherent risk is not the goal
				\item \textbf{C - Incorrect:} Can focus on reducing either frequency or impact
				\item \textbf{D - Incorrect:} Can focus on reducing either frequency or impact
			\end{itemize}
		\end{block}
	\end{frame}
	
	\begin{frame}{Q19: Best Control to Prevent Unauthorized Access}
		\begin{block}{Question}
			BEST control to prevent unauthorized access to information is user:
		\end{block}
		\begin{enumerate}[(A)]
			\item Accountability
			\item Authentication
			\item Identification
			\item Access rules
		\end{enumerate}
	\end{frame}
	
	\begin{frame}{Q19: Explanation}
		\begin{block}{Correct Answer: B}
			Authentication verifies the user's identity and the right to access information according to the access rules.
		\end{block}
		\begin{block}{Option Analysis}
			\begin{itemize}
				\item \textbf{A - Incorrect:} Accountability maps activity back to responsible party
				\item \textbf{B - Correct:} Authentication verifies identity
				\item \textbf{C - Incorrect:} Identification without authentication does not grant access
				\item \textbf{D - Incorrect:} Access rules without identification and authentication do not grant access
			\end{itemize}
		\end{block}
	\end{frame}
	
	\begin{frame}{Q20: Controlling Access to Sensitive Information}
		\begin{block}{Question}
			Which control BEST protects from unauthorized individuals gaining access to sensitive information?
		\end{block}
		\begin{enumerate}[(A)]
			\item Using a challenge response system
			\item Forcing periodic password changes
			\item Monitoring unsuccessful login attempts
			\item Providing access on a need-to-know basis
		\end{enumerate}
	\end{frame}
	
	\begin{frame}{Q20: Explanation}
		\begin{block}{Correct Answer: D}
			Physical or logical system access should be assigned on a need-to-know basis (legitimate business requirements) and in ways that incorporate least privilege and segregation of duties.
		\end{block}
		\begin{block}{Option Analysis}
			\begin{itemize}
				\item \textbf{A - Incorrect:} Challenge response doesn't address access design
				\item \textbf{B - Incorrect:} Password changes don't guarantee appropriate access
				\item \textbf{C - Incorrect:} Monitoring doesn't address appropriate access rights
				\item \textbf{D - Correct:} Need-to-know basis is fundamental
			\end{itemize}
		\end{block}
	\end{frame}
	
	\begin{frame}{Q21: Best Defense Against Phishing}
		\begin{block}{Question}
			Which defense is BEST to use against phishing attacks?
		\end{block}
		\begin{enumerate}[(A)]
			\item Intrusion detection system
			\item Spam filters
			\item End-user awareness
			\item Application hardening
		\end{enumerate}
	\end{frame}
	
	\begin{frame}{Q21: Explanation}
		\begin{block}{Correct Answer: C}
			Phishing attacks are a type of social engineering attack and are best defended by end-user awareness training.
		\end{block}
		\begin{block}{Option Analysis}
			\begin{itemize}
				\item \textbf{A - Incorrect:} IDS doesn't protect against phishing
				\item \textbf{B - Incorrect:} Spam filters are not as effective as user awareness
				\item \textbf{C - Correct:} User awareness is the best defense
				\item \textbf{D - Incorrect:} Application hardening doesn't protect against phishing
			\end{itemize}
		\end{block}
	\end{frame}
	
	\begin{frame}{Q22: Stakeholders in Risk Response Review}
		\begin{block}{Question}
			When responding to identified risk, MOST important stakeholders involved in reviewing risk response options are:
		\end{block}
		\begin{enumerate}[(A)]
			\item Information security managers
			\item Internal auditors
			\item Incident response team members
			\item Business managers
		\end{enumerate}
	\end{frame}
	
	\begin{frame}{Q22: Explanation}
		\begin{block}{Correct Answer: D}
			Business managers are accountable for managing the associated risk and will determine what actions to take based on the information provided by others, which may include collaboration with and support from IT security managers.
		\end{block}
		\begin{block}{Option Analysis}
			\begin{itemize}
				\item \textbf{A - Incorrect:} Security managers understand technical aspects but don't decide
				\item \textbf{B - Incorrect:} Risk response is not a function of internal audit
				\item \textbf{C - Incorrect:} Incident response team responds to incidents
				\item \textbf{D - Correct:} Business managers are accountable
			\end{itemize}
		\end{block}
	\end{frame}
	
	\begin{frame}{Q23: Designing Information System Controls}
		\begin{block}{Question}
			Which should be considered FIRST when designing information system controls?
		\end{block}
		\begin{enumerate}[(A)]
			\item Organizational strategic plan
			\item Existing IT environment
			\item Present IT budget
			\item IT strategic plan
		\end{enumerate}
	\end{frame}
	
	\begin{frame}{Q23: Explanation}
		\begin{block}{Correct Answer: A}
			Review of the enterprise's strategic plan is the first step in designing effective IS controls that would fit the enterprise's long-term plans.
		\end{block}
		\begin{block}{Option Analysis}
			\begin{itemize}
				\item \textbf{A - Correct:} Strategic plan is the starting point
				\item \textbf{B - Incorrect:} Useful but not the first task
				\item \textbf{C - Incorrect:} Budget is one component of strategic plan
				\item \textbf{D - Incorrect:} IT strategic plan exists to support enterprise's strategic plan
			\end{itemize}
		\end{block}
	\end{frame}
	
	\begin{frame}{Q24: Residual Risk Calculation Basis}
		\begin{block}{Question}
			Residual risk can be accurately calculated on the basis of:
		\end{block}
		\begin{enumerate}[(A)]
			\item Threats and vulnerabilities
			\item Inherent risk and control risk
			\item Compliance risk and reputation
			\item Risk governance and risk response
		\end{enumerate}
	\end{frame}
	
	\begin{frame}{Q24: Explanation}
		\begin{block}{Correct Answer: B}
			Inherent risk multiplied by control risk is the formula to calculate residual risk.
		\end{block}
		\begin{block}{Option Analysis}
			\begin{itemize}
				\item \textbf{A - Incorrect:} Threat/vulnerability are elements of inherent risk
				\item \textbf{B - Correct:} Inherent risk × control risk = residual risk
				\item \textbf{C - Incorrect:} Compliance risk is not a basis for calculation
				\item \textbf{D - Incorrect:} These are risk domains, not elements
			\end{itemize}
		\end{block}
	\end{frame}
	
	% ================ SECTION 4: DOMAIN 4 - RISK MONITORING ================
	
	\section{Domain 4: Risk and Control Monitoring}
	
	\begin{frame}{Q25: Maintaining KRIs}
		\begin{block}{Question}
			MOST important reason to maintain key risk indicators (KRIs) is that:
		\end{block}
		\begin{enumerate}[(A)]
			\item Complex metrics require fine-tuning
			\item Threats and vulnerabilities change over time
			\item Risk reports need to be timely
			\item They help to avoid risk
		\end{enumerate}
	\end{frame}
	
	\begin{frame}{Q25: Explanation}
		\begin{block}{Correct Answer: B}
			Threats and vulnerabilities change over time, and KRI maintenance ensures that KRIs continue to effectively capture these changes.
		\end{block}
		\begin{block}{Option Analysis}
			\begin{itemize}
				\item \textbf{A - Incorrect:} Fine-tuning is secondary to capturing changes
				\item \textbf{B - Correct:} Threats and vulnerabilities change over time
				\item \textbf{C - Incorrect:} Timeliness is a business requirement
				\item \textbf{D - Incorrect:} Risk avoidance is one risk response
			\end{itemize}
		\end{block}
	\end{frame}
	
	\begin{frame}{Q26: Measuring Operational Effectiveness}
		\begin{block}{Question}
			Which is BEST measure of operational effectiveness of risk management process capabilities?
		\end{block}
		\begin{enumerate}[(A)]
			\item Key performance indicators (KPIs)
			\item Key risk indicators (KRIs)
			\item Base practices
			\item Metric thresholds
		\end{enumerate}
	\end{frame}
	
	\begin{frame}{Q26: Explanation}
		\begin{block}{Correct Answer: A}
			Key performance indicators (KPIs) are assessment indicators that support judgment regarding the performance of a specific process.
		\end{block}
		\begin{block}{Option Analysis}
			\begin{itemize}
				\item \textbf{A - Correct:} KPIs measure process performance
				\item \textbf{B - Incorrect:} KRIs provide insights into potential risk
				\item \textbf{C - Incorrect:} Base practices need work products for evidence
				\item \textbf{D - Incorrect:} Thresholds are decision points
			\end{itemize}
		\end{block}
	\end{frame}
	
	\begin{frame}{Q27: Data Extraction Forecasting}
		\begin{block}{Question}
			During data extraction, total transactions per year was forecasted by multiplying monthly average by twelve. This is considered:
		\end{block}
		\begin{enumerate}[(A)]
			\item A controls total
			\item Simplistic and ineffective
			\item A duplicates test
			\item A reasonableness test
		\end{enumerate}
	\end{frame}
	
	\begin{frame}{Q27: Explanation}
		\begin{block}{Correct Answer: D}
			Reasonableness tests make certain assumptions about the information as the basis for more elaborate data validation tests.
		\end{block}
		\begin{block}{Option Analysis}
			\begin{itemize}
				\item \textbf{A - Incorrect:} Does not ensure all transactions extracted
				\item \textbf{B - Incorrect:} Reasonableness test is valid foundation
				\item \textbf{C - Incorrect:} Does not identify duplicate transactions
				\item \textbf{D - Correct:} This is a reasonableness test
			\end{itemize}
		\end{block}
	\end{frame}
	
	\begin{frame}{Q28: Testing Access Management Effectiveness}
		\begin{block}{Question}
			BEST test for confirming effectiveness of system access management process is to map:
		\end{block}
		\begin{enumerate}[(A)]
			\item Access requests to user accounts
			\item User accounts to access requests
			\item User accounts to HR records
			\item Vendor database to user accounts
		\end{enumerate}
	\end{frame}
	
	\begin{frame}{Q28: Explanation}
		\begin{block}{Correct Answer: B}
			Mapping user accounts to access requests confirms that all existing accounts have been approved.
		\end{block}
		\begin{block}{Option Analysis}
			\begin{itemize}
				\item \textbf{A - Incorrect:} Confirms requests processed but not unapproved accounts
				\item \textbf{B - Correct:} Confirms all accounts have been approved
				\item \textbf{C - Incorrect:} Confirms accounts tied to employees only
				\item \textbf{D - Incorrect:} Flawed because it doesn't consider unapproved accounts
			\end{itemize}
		\end{block}
	\end{frame}
	
	\begin{frame}{Q29: Firewall Compliance Assurance}
		\begin{block}{Question}
			Which provides BEST assurance that a firewall is configured in compliance with enterprise's security policy?
		\end{block}
		\begin{enumerate}[(A)]
			\item Reviewing actual procedures
			\item Interviewing firewall administrator
			\item Reviewing parameter settings
			\item Reviewing log file for recent attacks
		\end{enumerate}
	\end{frame}
	
	\begin{frame}{Q29: Explanation}
		\begin{block}{Correct Answer: C}
			A review of the parameter settings provides a good basis for comparison of the actual configuration to the security policy and reliable audit evidence documentation.
		\end{block}
		\begin{block}{Option Analysis}
			\begin{itemize}
				\item \textbf{A - Incorrect:} Procedures show how it should be managed, not compliance
				\item \textbf{B - Incorrect:} Interview provides overview, not confirmation
				\item \textbf{C - Correct:} Parameter settings show actual configuration
				\item \textbf{D - Incorrect:} Log review shows logging is enabled, not configuration
			\end{itemize}
		\end{block}
	\end{frame}
	
	\begin{frame}{Q30: Verifying Control Effectiveness}
		\begin{block}{Question}
			One way to verify control effectiveness is by determining:
		\end{block}
		\begin{enumerate}[(A)]
			\item Its reliability
			\item Whether it is preventive or detective
			\item Capability of providing notification of failure
			\item Test results of intended objectives
		\end{enumerate}
	\end{frame}
	
	\begin{frame}{Q30: Explanation}
		\begin{block}{Correct Answer: D}
			Control effectiveness requires a process to verify that the control process worked as intended and meets the intended control objectives.
		\end{block}
		\begin{block}{Option Analysis}
			\begin{itemize}
				\item \textbf{A - Incorrect:} Weak controls can be highly reliable
				\item \textbf{B - Incorrect:} Control type doesn't determine effectiveness
				\item \textbf{C - Incorrect:} Notification of failure doesn't determine strength
				\item \textbf{D - Correct:} Test results verify intended objectives
			\end{itemize}
		\end{block}
	\end{frame}
	
	\begin{frame}{Q31: Correlation Tools for Trend Analysis}
		\begin{block}{Question}
			Tools that correlate information from multiple systems to improve trend analysis are MOST likely applied to:
		\end{block}
		\begin{enumerate}[(A)]
			\item Transaction data
			\item Configuration settings
			\item System changes
			\item Process integrity
		\end{enumerate}
	\end{frame}
	
	\begin{frame}{Q31: Explanation}
		\begin{block}{Correct Answer: A}
			Transaction data tends to be difficult to analyze for trends across multiple systems because of sheer volume. Therefore, correlation engines are often used to analyze the data to see trends that otherwise might not be apparent.
		\end{block}
		\begin{block}{Option Analysis}
			\begin{itemize}
				\item \textbf{A - Correct:} Transaction data benefits from correlation
				\item \textbf{B - Incorrect:} Configuration settings are compared to predefined values
				\item \textbf{C - Incorrect:} System changes rarely benefit from trend analysis
				\item \textbf{D - Incorrect:} Process integrity not typically associated with trend analysis
			\end{itemize}
		\end{block}
	\end{frame}
	
	\begin{frame}{Q32: Ensuring Outsourced Compliance}
		\begin{block}{Question}
			Most effective way to ensure outsourced service providers comply with enterprise information security policy?
		\end{block}
		\begin{enumerate}[(A)]
			\item Periodic audits
			\item Security awareness training
			\item Penetration testing
			\item Service level monitoring
		\end{enumerate}
	\end{frame}
	
	\begin{frame}{Q32: Explanation}
		\begin{block}{Correct Answer: A}
			Regular audits can identify gaps in information security compliance.
		\end{block}
		\begin{block}{Option Analysis}
			\begin{itemize}
				\item \textbf{A - Correct:} Audits ensure compliance
				\item \textbf{B - Incorrect:} Training increases awareness but less effective
				\item \textbf{C - Incorrect:} Penetration testing identifies vulnerabilities
				\item \textbf{D - Incorrect:} Service monitoring identifies operational issues
			\end{itemize}
		\end{block}
	\end{frame}
	
	% ================ ADDITIONAL DOMAIN 1 QUESTIONS ================
	
	\section{Domain 1: Additional Key Concepts}
	
	\begin{frame}{Q33: Risk Appetite and Tolerance Definition}
		\begin{block}{Question}
			Risk appetite is defined as:
		\end{block}
		\begin{enumerate}[(A)]
			\item The amount of risk an enterprise can tolerate without existence being questioned
			\item The amount of risk an enterprise is willing to accept in pursuit of its mission
			\item The acceptable level of variation management is willing to allow
			\item The maximum loss an enterprise can absorb
		\end{enumerate}
	\end{frame}
	
	\begin{frame}{Q33: Explanation}
		\begin{block}{Correct Answer: B}
			Risk appetite is defined as the amount of risk, on a broad level, that an entity is willing to accept in pursuit of its mission.
		\end{block}
		\begin{block}{Option Analysis}
			\begin{itemize}
				\item \textbf{A - Incorrect:} This is risk capacity
				\item \textbf{B - Correct:} This is the definition of risk appetite
				\item \textbf{C - Incorrect:} This is risk tolerance
				\item \textbf{D - Incorrect:} This is risk capacity
			\end{itemize}
		\end{block}
	\end{frame}
	
	\begin{frame}{Q34: Risk Tolerance Example}
		\begin{block}{Question}
			Risk tolerance is best described as:
		\end{block}
		\begin{enumerate}[(A)]
			\item The maximum risk an enterprise can take
			\item The acceptable level of deviation from risk appetite
			\item The complete elimination of risk
			\item The transfer of risk to a third party
		\end{enumerate}
	\end{frame}
	
	\begin{frame}{Q34: Explanation}
		\begin{block}{Correct Answer: B}
			Risk tolerance is defined as the acceptable level of variation that management is willing to allow for any particular risk as the enterprise pursues its objectives.
		\end{block}
		\begin{block}{Option Analysis}
			\begin{itemize}
				\item \textbf{A - Incorrect:} This is risk capacity
				\item \textbf{B - Correct:} Tolerance is acceptable deviation
				\item \textbf{C - Incorrect:} Risk cannot be completely eliminated
				\item \textbf{D - Incorrect:} This is risk transfer
			\end{itemize}
		\end{block}
	\end{frame}
	
	\begin{frame}{Q35: Risk Culture Symptoms}
		\begin{block}{Question}
			Which is a symptom of an inadequate or problematic risk culture?
		\end{block}
		\begin{enumerate}[(A)]
			\item Open communication about risk
			\item Clear understanding of risk appetite
			\item Overconfidence in ability to manage risk
			\item Regular risk reporting to management
		\end{enumerate}
	\end{frame}
	
	\begin{frame}{Q35: Explanation}
		\begin{block}{Correct Answer: C}
			Overconfidence in ability to manage risk is a symptom of problematic risk culture.
		\end{block}
		\begin{block}{Option Analysis}
			\begin{itemize}
				\item \textbf{A - Incorrect:} Open communication is healthy
				\item \textbf{B - Incorrect:} Clear understanding is positive
				\item \textbf{C - Correct:} Overconfidence indicates problematic culture
				\item \textbf{D - Incorrect:} Regular reporting is good practice
			\end{itemize}
		\end{block}
	\end{frame}
	
	\begin{frame}{Q36: Confidentiality Breach}
		\begin{block}{Question}
			Which is an example of a breach of confidentiality?
		\end{block}
		\begin{enumerate}[(A)]
			\item Data being modified by unauthorized user
			\item Data being unavailable when needed
			\item Data being disclosed to unauthorized recipient
			\item Data being destroyed
		\end{enumerate}
	\end{frame}
	
	\begin{frame}{Q36: Explanation}
		\begin{block}{Correct Answer: C}
			A breach of confidentiality means the improper disclosure of information to a recipient not authorized access by the information owner.
		\end{block}
		\begin{block}{Option Analysis}
			\begin{itemize}
				\item \textbf{A - Incorrect:} This is a breach of integrity
				\item \textbf{B - Incorrect:} This is a breach of availability
				\item \textbf{C - Correct:} Unauthorized disclosure is confidentiality breach
				\item \textbf{D - Incorrect:} This is a breach of integrity/availability
			\end{itemize}
		\end{block}
	\end{frame}
	
	\begin{frame}{Q37: CIA Triad Trade-off}
		\begin{block}{Question}
			Increasing confidentiality typically:
		\end{block}
		\begin{enumerate}[(A)]
			\item Increases availability
			\item Decreases processing time
			\item Increases processing time and may reduce availability
			\item Has no effect on availability
		\end{enumerate}
	\end{frame}
	
	\begin{frame}{Q37: Explanation}
		\begin{block}{Correct Answer: C}
			Increasing confidentiality typically increases processing time, which reduces availability. There are trade-offs between the three principles.
		\end{block}
		\begin{block}{Option Analysis}
			\begin{itemize}
				\item \textbf{A - Incorrect:} Usually decreases availability
				\item \textbf{B - Incorrect:} Usually increases processing time
				\item \textbf{C - Correct:} Trade-off between confidentiality and availability
				\item \textbf{D - Incorrect:} There is typically an effect
			\end{itemize}
		\end{block}
	\end{frame}
	
	% ================ ADDITIONAL DOMAIN 2 QUESTIONS ================
	
	\section{Domain 2: Additional Assessment Concepts}
	
	\begin{frame}{Q38: OCTAVE Methodology}
		\begin{block}{Question}
			OCTAVE focuses primarily on:
		\end{block}
		\begin{enumerate}[(A)]
			\item Technical vulnerabilities only
			\item Critical assets and risk to those assets
			\item Financial risk quantification
			\item Compliance requirements
		\end{enumerate}
	\end{frame}
	
	\begin{frame}{Q38: Explanation}
		\begin{block}{Correct Answer: B}
			OCTAVE focuses on critical assets and the risk to those assets using a comprehensive, systematic, context-driven and self-directed evaluation approach.
		\end{block}
		\begin{block}{Option Analysis}
			\begin{itemize}
				\item \textbf{A - Incorrect:} Not focused only on technical vulnerabilities
				\item \textbf{B - Correct:} Focuses on critical assets
				\item \textbf{C - Incorrect:} Not financial quantification
				\item \textbf{D - Incorrect:} Not primarily compliance
			\end{itemize}
		\end{block}
	\end{frame}
	
	\begin{frame}{Q39: Phases of OCTAVE}
		\begin{block}{Question}
			Phase 1 of OCTAVE involves:
		\end{block}
		\begin{enumerate}[(A)]
			\item Identifying infrastructure vulnerabilities
			\item Building asset-based threat profiles
			\item Developing security strategy
			\item Implementing controls
		\end{enumerate}
	\end{frame}
	
	\begin{frame}{Q39: Explanation}
		\begin{block}{Correct Answer: B}
			Phase 1 of OCTAVE is to build asset-based threat profiles through organizational evaluation.
		\end{block}
		\begin{block}{Option Analysis}
			\begin{itemize}
				\item \textbf{A - Incorrect:} This is Phase 2
				\item \textbf{B - Correct:} Phase 1 builds threat profiles
				\item \textbf{C - Incorrect:} This is Phase 3
				\item \textbf{D - Incorrect:} This follows OCTAVE
			\end{itemize}
		\end{block}
	\end{frame}
	
	\begin{frame}{Q40: Fault Tree Analysis}
		\begin{block}{Question}
			Fault tree analysis is best described as:
		\end{block}
		\begin{enumerate}[(A)]
			\item Forward-looking, bottom-up model
			\item Top-down analysis starting with an event
			\item Statistical inference method
			\item Delphi consensus technique
		\end{enumerate}
	\end{frame}
	
	\begin{frame}{Q40: Explanation}
		\begin{block}{Correct Answer: B}
			A fault tree analysis starts with an event and examines possible means for the event to occur (top-down) and displays these results in a logical tree diagram.
		\end{block}
		\begin{block}{Option Analysis}
			\begin{itemize}
				\item \textbf{A - Incorrect:} This is event tree analysis
				\item \textbf{B - Correct:} Top-down analysis starting with event
				\item \textbf{C - Incorrect:} This is Bayesian analysis
				\item \textbf{D - Incorrect:} This is Delphi method
			\end{itemize}
		\end{block}
	\end{frame}
	
	\begin{frame}{Q41: Monte Carlo Analysis}
		\begin{block}{Question}
			Monte Carlo analysis is used to:
		\end{block}
		\begin{enumerate}[(A)]
			\item Eliminate risk entirely
			\item Establish aggregate variation in a system
			\item Identify single points of failure
			\item Document security policies
		\end{enumerate}
	\end{frame}
	
	\begin{frame}{Q41: Explanation}
		\begin{block}{Correct Answer: B}
			Monte Carlo simulation is used to establish the aggregate variation in a system resulting from variations in the system, for a number of inputs, where each input has a defined distribution.
		\end{block}
		\begin{block}{Option Analysis}
			\begin{itemize}
				\item \textbf{A - Incorrect:} Cannot eliminate risk entirely
				\item \textbf{B - Correct:} Establishes aggregate variation
				\item \textbf{C - Incorrect:} Not for identifying SPOFs
				\item \textbf{D - Incorrect:} Not for documentation
			\end{itemize}
		\end{block}
	\end{frame}
	
	\begin{frame}{Q42: Delphi Method}
		\begin{block}{Question}
			Delphi method leverages:
		\end{block}
		\begin{enumerate}[(A)]
			\item Expert opinion through multiple questionnaire rounds
			\item Statistical analysis of historical data
			\item Automated vulnerability scanning
			\item Document review
		\end{enumerate}
	\end{frame}
	
	\begin{frame}{Q42: Explanation}
		\begin{block}{Correct Answer: A}
			The Delphi method leverages expert opinion received using two or more rounds of questionnaires. After each round, results are summarized and communicated by a facilitator.
		\end{block}
		\begin{block}{Option Analysis}
			\begin{itemize}
				\item \textbf{A - Correct:} Expert opinion through multiple rounds
				\item \textbf{B - Incorrect:} Not statistical analysis
				\item \textbf{C - Incorrect:} Not automated scanning
				\item \textbf{D - Incorrect:} Not document review
			\end{itemize}
		\end{block}
	\end{frame}
	
	% ================ ADDITIONAL DOMAIN 3 QUESTIONS ================
	
	\section{Domain 3: Additional Response Concepts}
	
	\begin{frame}{Q43: Cost-Benefit Analysis Purpose}
		\begin{block}{Question}
			Cost-benefit analysis is used to:
		\end{block}
		\begin{enumerate}[(A)]
			\item Justify expense of control implementation
			\item Calculate annual loss expectancy
			\item Determine risk appetite
			\item Identify threats
		\end{enumerate}
	\end{frame}
	
	\begin{frame}{Q43: Explanation}
		\begin{block}{Correct Answer: A}
			Cost-benefit analysis is used to justify the expense associated with the implementation of controls. The expenditure on a control cannot be justified if the benefit realized is less than the cost.
		\end{block}
		\begin{block}{Option Analysis}
			\begin{itemize}
				\item \textbf{A - Correct:} Justifies control expense
				\item \textbf{B - Incorrect:} ALE is calculated differently
				\item \textbf{C - Incorrect:} Risk appetite is management decision
				\item \textbf{D - Incorrect:} Threats are identified through risk identification
			\end{itemize}
		\end{block}
	\end{frame}
	
	\begin{frame}{Q44: Risk Transfer Example}
		\begin{block}{Question}
			Risk transfer is best exemplified by:
		\end{block}
		\begin{enumerate}[(A)]
			\item Purchasing insurance
			\item Implementing a firewall
			\item Conducting user training
			\item Creating a policy
		\end{enumerate}
	\end{frame}
	
	\begin{frame}{Q44: Explanation}
		\begin{block}{Correct Answer: A}
			Risk transfer is a decision to reduce loss by having another organization incur the cost. The most common example of risk transfer is the purchasing of insurance.
		\end{block}
		\begin{block}{Option Analysis}
			\begin{itemize}
				\item \textbf{A - Correct:} Insurance transfers risk
				\item \textbf{B - Incorrect:} Firewall is mitigation
				\item \textbf{C - Incorrect:} Training is mitigation/awareness
				\item \textbf{D - Incorrect:} Policy is administrative control
			\end{itemize}
		\end{block}
	\end{frame}
	
	\begin{frame}{Q45: When to Accept Risk}
		\begin{block}{Question}
			Risk should be accepted when:
		\end{block}
		\begin{enumerate}[(A)]
			\item It falls within organizational risk appetite
			\item It exceeds risk capacity
			\item Controls are available
			\item The cost of controls is low
		\end{enumerate}
	\end{frame}
	
	\begin{frame}{Q45: Explanation}
		\begin{block}{Correct Answer: A}
			Risk that falls within the organizational risk appetite should be accepted. The decision to accept risk is made according to the risk appetite and risk tolerance set by senior management.
		\end{block}
		\begin{block}{Option Analysis}
			\begin{itemize}
				\item \textbf{A - Correct:} Within risk appetite
				\item \textbf{B - Incorrect:} Cannot exceed risk capacity
				\item \textbf{C - Incorrect:} If controls available, consider mitigation
				\item \textbf{D - Incorrect:} Low cost may justify mitigation
			\end{itemize}
		\end{block}
	\end{frame}
	
	\begin{frame}{Q46: Risk Avoidance}
		\begin{block}{Question}
			Risk avoidance means:
		\end{block}
		\begin{enumerate}[(A)]
			\item Transferring risk to insurance
			\item Exiting activities that give rise to risk
			\item Reducing risk through controls
			\item Accepting risk without action
		\end{enumerate}
	\end{frame}
	
	\begin{frame}{Q46: Explanation}
		\begin{block}{Correct Answer: B}
			Risk avoidance means exiting the activities or conditions that give rise to risk. It is the choice that remains when no other response is adequate.
		\end{block}
		\begin{block}{Option Analysis}
			\begin{itemize}
				\item \textbf{A - Incorrect:} This is risk transfer
				\item \textbf{B - Correct:} Exiting risk-causing activities
				\item \textbf{C - Incorrect:} This is risk mitigation
				\item \textbf{D - Incorrect:} This is risk acceptance
			\end{itemize}
		\end{block}
	\end{frame}
	
	\begin{frame}{Q47: Symmetric Encryption}
		\begin{block}{Question}
			Symmetric encryption uses:
		\end{block}
		\begin{enumerate}[(A)]
			\item Two mathematically related keys
			\item The same secret key for encryption and decryption
			\item Public and private key pairs
			\item Digital certificates
		\end{enumerate}
	\end{frame}
	
	\begin{frame}{Q47: Explanation}
		\begin{block}{Correct Answer: B}
			Symmetric key cryptography occurs when an encryption algorithm uses the same secret key to encrypt the plaintext into ciphertext and to decrypt the ciphertext back into plaintext.
		\end{block}
		\begin{block}{Option Analysis}
			\begin{itemize}
				\item \textbf{A - Incorrect:} This is asymmetric encryption
				\item \textbf{B - Correct:} Same key for both operations
				\item \textbf{C - Incorrect:} This is asymmetric encryption
				\item \textbf{D - Incorrect:} Certificates are for identity verification
			\end{itemize}
		\end{block}
	\end{frame}
	
	\begin{frame}{Q48: Digital Signature}
		\begin{block}{Question}
			A digital signature provides:
		\end{block}
		\begin{enumerate}[(A)]
			\item Confidentiality only
			\item Integrity and nonrepudiation
			\item Availability only
			\item Authentication only
		\end{enumerate}
	\end{frame}
	
	\begin{frame}{Q48: Explanation}
		\begin{block}{Correct Answer: B}
			A digital signature combines a hash function with asymmetric encryption to verify the author's identity. It provides integrity (message not changed) and nonrepudiation (sender cannot deny sending).
		\end{block}
		\begin{block}{Option Analysis}
			\begin{itemize}
				\item \textbf{A - Incorrect:} Does not provide confidentiality on its own
				\item \textbf{B - Correct:} Provides integrity and nonrepudiation
				\item \textbf{C - Incorrect:} Not related to availability
				\item \textbf{D - Incorrect:} Provides more than authentication
			\end{itemize}
		\end{block}
	\end{frame}
	
	\begin{frame}{Q49: Asymmetric Encryption Advantage}
		\begin{block}{Question}
			Asymmetric encryption's main advantage over symmetric is:
		\end{block}
		\begin{enumerate}[(A)]
			\item Faster processing speed
			\item Easier key management for secure exchange
			\item Lower computational requirements
			\item Simpler implementation
		\end{enumerate}
	\end{frame}
	
	\begin{frame}{Q49: Explanation}
		\begin{block}{Correct Answer: B}
			Asymmetric algorithms overcome the weakness of scalability of symmetric key cryptography by allowing secure key exchange. The public key can be freely distributed, making key management easier.
		\end{block}
		\begin{block}{Option Analysis}
			\begin{itemize}
				\item \textbf{A - Incorrect:} Asymmetric is slower
				\item \textbf{B - Correct:} Easier key management
				\item \textbf{C - Incorrect:} Higher computational requirements
				\item \textbf{D - Incorrect:} More complex implementation
			\end{itemize}
		\end{block}
	\end{frame}
	
	% ================ ADDITIONAL DOMAIN 4 QUESTIONS ================
	
	\section{Domain 4: Additional Monitoring Concepts}
	
	\begin{frame}{Q50: KRI Sensitivity}
		\begin{block}{Question}
			A KRI that generates too many false positives should be adjusted for:
		\end{block}
		\begin{enumerate}[(A)]
			\item Sensitivity
			\item Timeliness
			\item Frequency
			\item Reliability
		\end{enumerate}
	\end{frame}
	
	\begin{frame}{Q50: Explanation}
		\begin{block}{Correct Answer: A}
			KRIs that report on data points that cannot be controlled or are not alerting at the correct time should be adjusted (optimized) to be more precise, relevant or accurate. Sensitivity should be adjusted to reduce false positives.
		\end{block}
		\begin{block}{Option Analysis}
			\begin{itemize}
				\item \textbf{A - Correct:} Sensitivity affects false positives
				\item \textbf{B - Incorrect:} Timing is about when alerts occur
				\item \textbf{C - Incorrect:} Frequency is how often measured
				\item \textbf{D - Incorrect:} Reliability is correlation with risk
			\end{itemize}
		\end{block}
	\end{frame}
	
	\begin{frame}{Q51: KPI and KRI Relationship}
		\begin{block}{Question}
			KPIs and KRIs are often used together because:
		\end{block}
		\begin{enumerate}[(A)]
			\item They measure the same thing
			\item KPIs identify underperformance; KRIs provide early warnings
			\item KRIs replace KPIs in mature organizations
			\item They are independent and not related
		\end{enumerate}
	\end{frame}
	
	\begin{frame}{Q51: Explanation}
		\begin{block}{Correct Answer: B}
			KPIs help to identify underperforming aspects of the organization, while KRIs provide early warnings of increased risk. Failing to meet performance goals is a risk, so the two can be used together.
		\end{block}
		\begin{block}{Option Analysis}
			\begin{itemize}
				\item \textbf{A - Incorrect:} They measure different things
				\item \textbf{B - Correct:} Complementary purposes
				\item \textbf{C - Incorrect:} KRIs don't replace KPIs
				\item \textbf{D - Incorrect:} They are related
			\end{itemize}
		\end{block}
	\end{frame}
	
	\begin{frame}{Q52: KRI Maintenance Frequency}
		\begin{block}{Question}
			KRIs should be evaluated:
		\end{block}
		\begin{enumerate}[(A)]
			\item Annually
			\item Only when incidents occur
			\item On a regular basis
			\item Never after implementation
		\end{enumerate}
	\end{frame}
	
	\begin{frame}{Q52: Explanation}
		\begin{block}{Correct Answer: C}
			Because the organization's internal and external environments are constantly changing, the set of KRIs needs to be evaluated on a regular basis to verify that each KRI remains properly related to the risk appetite and tolerance levels.
		\end{block}
		\begin{block}{Option Analysis}
			\begin{itemize}
				\item \textbf{A - Incorrect:} May need more frequent evaluation
				\item \textbf{B - Incorrect:} Should be proactive, not reactive
				\item \textbf{C - Correct:} Regular evaluation is needed
				\item \textbf{D - Incorrect:} Must be maintained over time
			\end{itemize}
		\end{block}
	\end{frame}
	
	\begin{frame}{Q53: SIEM Primary Function}
		\begin{block}{Question}
			SIEM systems primarily:
		\end{block}
		\begin{enumerate}[(A)]
			\item Correlate data from multiple sources
			\item Store backup data
			\item Manage user identities
			\item Encrypt network traffic
		\end{enumerate}
	\end{frame}
	
	\begin{frame}{Q53: Explanation}
		\begin{block}{Correct Answer: A}
			SIEM systems are integrated data correlation tools that capture data from multiple sources and analyze system, application and network activity for possible security events.
		\end{block}
		\begin{block}{Option Analysis}
			\begin{itemize}
				\item \textbf{A - Correct:} Correlation of multiple data sources
				\item \textbf{B - Incorrect:} Not for backup storage
				\item \textbf{C - Incorrect:} Identity management is separate
				\item \textbf{D - Incorrect:} Not for traffic encryption
			\end{itemize}
		\end{block}
	\end{frame}
	
	\begin{frame}{Q54: Log Review Challenges}
		\begin{block}{Question}
			Common challenge with log review is:
		\end{block}
		\begin{enumerate}[(A)]
			\item Too much data to analyze effectively
			\item Not enough data being collected
			\item Logs being too expensive
			\item Logs not being required
		\end{enumerate}
	\end{frame}
	
	\begin{frame}{Q54: Explanation}
		\begin{block}{Correct Answer: A}
			Common challenges relating to the effective use of logs include having too much data, difficulty in searching for relevant information, improper configuration, and modification or deletion of data before it is read.
		\end{block}
		\begin{block}{Option Analysis}
			\begin{itemize}
				\item \textbf{A - Correct:} Data volume is a challenge
				\item \textbf{B - Incorrect:} Often there is too much data
				\item \textbf{C - Incorrect:} Cost is not the primary challenge
				\item \textbf{D - Incorrect:} Logs are typically required
			\end{itemize}
		\end{block}
	\end{frame}
	
	\begin{frame}{Q55: Control Monitoring Purpose}
		\begin{block}{Question}
			Purpose of control monitoring is to verify:
		\end{block}
		\begin{enumerate}[(A)]
			\item Whether the control is working
			\item Whether the control is effectively addressing the risk
			\item Whether staff is following procedures
			\item Whether logs are being reviewed
		\end{enumerate}
	\end{frame}
	
	\begin{frame}{Q55: Explanation}
		\begin{block}{Correct Answer: B}
			Because the purpose of a control is to mitigate a risk, the purpose of control monitoring is to verify whether the control is effectively addressing the risk, not to see whether the control itself is working.
		\end{block}
		\begin{block}{Option Analysis}
			\begin{itemize}
				\item \textbf{A - Incorrect:} This is just operational status
				\item \textbf{B - Correct:} Verify effectiveness against risk
				\item \textbf{C - Incorrect:} This is compliance monitoring
				\item \textbf{D - Incorrect:} This is log review, not control monitoring
			\end{itemize}
		\end{block}
	\end{frame}
	
	% ================ FINAL SLIDES ================
	
	\begin{frame}{CRISC Key Takeaways from Review Manual}
		\begin{block}{Essential Concepts}
			\begin{itemize}
				\item Risk appetite = amount of risk willing to accept in pursuit of mission
				\item Risk tolerance = acceptable level of deviation from appetite
				\item Risk capacity = maximum risk before existence is questioned
				\item C MM: 5 levels from Incomplete to Optimized
				\item CIA: Confidentiality, Integrity, Availability
				\item Digital signatures provide integrity and nonrepudiation
			\end{itemize}
		\end{block}
	\end{frame}
	
	\begin{frame}{CRISC Key Takeaways (continued)}
		\begin{block}{More Essential Concepts}
			\begin{itemize}
				\item Risk assessment = identify and evaluate risk and potential effects
				\item Risk response options: Accept, Mitigate, Transfer, Avoid
				\item Residual risk = Inherent risk - Controls
				\item KRIs predict risk; KPIs measure performance
				\item 3 LoD: Operational management, Risk oversight, Audit
			\end{itemize}
		\end{block}
	\end{frame}
	
	\begin{frame}{Final Advice}
		\begin{block}{Exam Preparation}
			\centering
			\vspace{0.3cm}
			\textbf{Remember:}
			\begin{itemize}
				\item Understand the \textit{why} behind each concept
				\item Know the difference between risk appetite, tolerance, and capacity
				\item Understand the risk management life cycle (Identify, Assess, Respond, Monitor)
				\item Practice scenario-based questions
				\item Use the ISACA mindset - choose the MOST appropriate answer
			\end{itemize}
			\vspace{0.5cm}
			\textbf{Good Luck on Your CRISC Exam!}
		\end{block}
	\end{frame}
	
\end{document}